\section{Scalar normalization on a universal cover}
\label{sec:scalar-normalization}

Starting from the homogeneous ray representatives $b_r\in A_0$ and their
structure constants $\kappa(r,s)\in k^*$ in \eqref{eq:projective-scalars}, we
wish to find nonzero scalars $a_r\in k^*$ such that the rescaled elements
$a_r b_r$ multiply with structure constants equal to $1$, that is,
\[
                   a_ra_s\kappa(r,s)=a_{rs}\qquad\text{whenever }rs\ne0.
\]
Because distinct pairs of rays can have the same nonzero product, these
equations are coupled across all corners of $A_0$. Following Bongartz and
Gabriel \cite[Section~4]{BongartzGabriel} and Bretscher and Gabriel
\cite[Sections~2.6 and~3.1--3.2]{BretscherGabriel}, we solve them by
constructing a two-dimensional cell complex $K$ from the Auslander--Reiten
quiver $\Gamma_A$ and its meshes, and lifting the mesh category $\C$ to the
universal cover $\widetilde K$. On the universal cover, the category of
lifted projectives is directed and thin, allowing us to construct a
multiplicative basis one vertex at a time. Because every multiplicative
scalar weight on the cover is a ratio of object values, and the fundamental
group $G=\pi_1(K)$ is free (Lemma~\ref{lem:mesh-walk-free}), we can then
rescale the lifted basis to make it $G$-invariant and descend it to $A_0$.

\subsection{A cell complex of meshes and its periodic components}
\label{subsec:mesh-walk-cover}

Since a basic algebra is a direct product of connected blocks, we may work
with one connected block at a time and assume throughout this section that
$A$ is connected; then $\Gamma_A$ is connected by
\cite[Theorem~VII.2.1]{ARS}. Write $\Gamma=\Gamma_A$, and identify $\C$ with
$\ind(\modu A_0)$ via Proposition~\ref{prop:projective-evaluation}. To form
the cell complex $K$ of \cite[Section~4.1]{BongartzGabriel}, start with the
underlying undirected graph of $\Gamma$, adjoin a translation edge
$t_X:\tau X\to X$ for every nonprojective vertex $X$, and for every arrow
$a:Y\to X$ with $X$ nonprojective, attach a two-dimensional triangle cell
whose boundary is the closed walk $(\sigma a)a t_X^{-1}$ from $\tau X$
through $Y$ and $X$ back to $\tau X$. Here a \emph{walk} in $K$ means a path
of edges traversed in either orientation; each triangle cell identifies the
two-edge walk $(\sigma a)a$ up to homotopy with the translation edge $t_X$
(rather than equating the individual summands of the mesh relation in the
linear category).

Let $\Gamma_{\mathrm{per}}$ denote the full subquiver of $\Gamma$ on the
vertices that are periodic under $\tau$, and call its connected components
the \emph{periodic components} of $\Gamma$. Because $\tau$ and $\tau^{-1}$
are everywhere defined on $\Gamma_{\mathrm{per}}$ and the mate bijection
$\sigma$ restricts to arrows between periodic vertices, $\Gamma_{\mathrm{per}}$
is a stable translation quiver. By \eqref{eq:mesh-object-quotient}, its
scalar mesh category is the object quotient of $\C$ obtained by killing all
nonperiodic vertices, and in particular is finite-dimensional. Riedtmann's
structure theorem \cite[Sections~1.5--1.6 and~2.1]{Riedtmann1980} describes
such stable components as quotients of repetition quivers $\Z\Delta$, where
for an oriented tree $\Delta$ the quiver $\Z\Delta$ has vertices
$(x,j)\in\Delta_0\times\Z$, arrows $(x,j)\to(y,j)$ and $(y,j)\to(x,j+1)$ for
each arrow $x\to y$ of $\Delta$, and translation $\tau(x,j)=(x,j-1)$. We
prove the following self-contained form in Appendix~\ref{app:periodic-trees}.

\begin{proposition}
\label{prop:periodic-tree-model}
Let $Q$ be a finite stable translation quiver with at most one arrow between
each ordered pair of vertices. Fix a coefficient system $c$ and suppose that
the scalar mesh category $C(Q,c)$ of Definition~\ref{def:scalar-mesh-category}
is finite-dimensional.
\begin{enumerate}[label={\upshape(\arabic*)}]
\item Every connected component of $Q$ containing at least one arrow is
isomorphic to $\Z\Delta/\Theta$, where $\Delta$ is a finite simply-laced
Dynkin tree and $\Theta$ is an infinite cyclic group of automorphisms of
$\Z\Delta$ acting freely on vertices and commuting with $\tau$.
\item If $\Delta$ has type $D_n$, $E_6$, $E_7$, $E_8$, or $A_{2n+1}$, root
$\Delta$ at its unique branch vertex (for types $D$ and $E$) or at its middle
vertex (for type $A_{2n+1}$), and orient all edges of $\Delta$ away from the
root. Then $\Theta$ is generated by $\tau^s\rho$ for some integer $s\geq1$,
where $\rho$ is the automorphism $(v,j)\mapsto(\rho v,j)$ of $\Z\Delta$
induced by a root-fixing automorphism $\rho$ of $\Delta$. If $\Delta=A_{2n}$,
then $\Theta$ is generated by $\eta^r$ for some integer $r\geq1$, where
$\eta$ is an automorphism of $\Z\Delta$ satisfying $\eta^2=\tau^{-1}$ that
interchanges the two halves of $\Delta$.
\end{enumerate}
\end{proposition}

Geometrically, the mesh triangles lying over each edge of $\Delta$ form an
infinite strip in $\Z\Delta$, so the subcomplex of $K$ on a periodic
component is the mapping torus of the tree $\Delta$ under the monodromy
automorphism generating $\Theta$. Using this description, we can deformation-retract
each periodic component onto a circle and $K$ itself onto a finite graph.

\begin{lemma}[{\cite[Theorem~4.2 and its corollary]{BongartzGabriel}}]
\label{lem:mesh-walk-free}
Let $K$ be the mesh cell complex of the connected Auslander--Reiten quiver
$\Gamma_A$ defined in Subsection~\ref{subsec:mesh-walk-cover}. Then $K$ is
homotopy equivalent to a finite graph, and consequently its fundamental
group $G=\pi_1(K)$ is free.
\end{lemma}
\begin{proof}
The partial mate bijection $\sigma$ partitions the arrows of $\Gamma_A$ into
finite chains and finite cycles. Since $\sigma^2=\tau$, an arrow lies in a
cycle of $\sigma$ if and only if both of its endpoints are $\tau$-periodic;
thus the cyclic $\sigma$-orbits are precisely the arrows of
$\Gamma_{\mathrm{per}}$. At the end of any finite $\sigma$-chain of length at
least two, the terminal arrow belongs to only one mesh triangle of $K$. We
may therefore collapse that free arrow together with its unique triangle, and
repeat the collapse inductively along each finite $\sigma$-chain until only a
single quiver edge per chain and the translation edges $t_X$ remain outside
the periodic subcomplexes.

Next, consider a periodic component $\Z\Delta/\Theta$, modeled as in
Proposition~\ref{prop:periodic-tree-model}. When $\Delta$ has a central root
vertex fixed by $\rho$, we inductively prune the outermost leaves of $\Delta$
in complete $\rho$-orbits. Specifically, for a leaf $v$ of $\Delta$ with
parent $u$, the arrows between their $\tau$-orbits in $\Z\Delta$ are
\[
 a_j:u_j\longrightarrow v_j,\qquad
 b_j:v_j\longrightarrow u_{j+1},
\]
and the incident mesh triangles have boundaries
$a_jb_jt_{u_{j+1}}^{-1}$ and $b_ja_{j+1}t_{v_{j+1}}^{-1}$. Because $v$ is a
leaf, the translation edge $t_{v_{j+1}}$ belongs only to the second triangle;
collapsing $t_{v_{j+1}}$ across that triangle leaves the edge $b_j$ incident
only to the first triangle, and collapsing $b_j$ across the first triangle
leaves only the edges $a_j:u_j\to v_j$, which then contract onto the parent
vertices $u_j$. Performing these elementary collapses simultaneously across
a $\rho$-orbit of leaves is compatible with the $\Theta$-action, and
repeating the pruning down to the root contracts the periodic component onto
the circle formed by the translation edges over the root.

When $\Delta=A_{2n}$, its center is the middle edge rather than a vertex.
Pruning all other edges as above reduces the component to the periodic strip
over the central edge. Subdividing this strip along its midline into two
half-strips gives a deformation retraction onto the midline that is
equivariant with respect to both $\tau$ and the involution $\eta$ exchanging
the two endpoints of the central edge; hence the quotient again contracts
onto a circle. Likewise, any periodic component consisting of a single
$\tau$-orbit with no arrows is already a circle of translation edges $t_X$.

Finally, after contracting all periodic subcomplexes to circles and
collapsing all nonperiodic triangles as in the first paragraph, only
one-dimensional cells remain, so $K$ deformation-retracts onto a finite
graph, whose fundamental group is free.
\end{proof}

Reversing the elementary triangle collapses of Lemma~\ref{lem:mesh-walk-free}
also shows that any closed edge weight on the periodic subcomplexes extends
to all of $K$. Here an \emph{edge weight} on $K$ with values in an abelian
group assigns opposite values to the two orientations of each edge, and it
is \emph{closed} if the sum of its values around the boundary of every
triangle cell is zero.

\begin{lemma}
\label{lem:periodic-weight-extension}
Every closed edge weight on the union of the periodic subcomplexes of $K$
extends to a closed edge weight on $K$.
\end{lemma}
\begin{proof}
As in the first paragraph of the proof of Lemma~\ref{lem:mesh-walk-free}, we
may collapse all triangles outside the periodic subcomplexes one by one along
finite $\sigma$-chains. Assign weight zero to all nonperiodic edges that
remain after these collapses, and then restore the collapsed triangles in
reverse order: since each restored triangle has a unique free edge that was
removed at that step, its weight is uniquely determined by requiring the
boundary sum around the triangle to vanish. This extends the given periodic
weight to a closed edge weight on all of $K$.
\end{proof}

As an immediate topological consequence of Lemma~\ref{lem:mesh-walk-free}, we
can normalize all scalar coefficients in the mesh relations of $k(\Gamma_A)$
to $1$.

\begin{lemma}
\label{lem:unit-mesh-coefficients}
In every characteristic, any system of nonzero mesh coefficients $c_{X,a}$ on
$\Gamma_A$ can be normalized to $1$ by rescaling the arrows and the mesh
relations.
\end{lemma}
\begin{proof}
Working on one connected component of $\Gamma_A$ at a time, view the family
of nonzero scalars $c_{X,a}\in k^*$ as a multiplicative cellular $2$-cochain
on $K$ assigning $c_{X,a}$ to the triangle cell with boundary
$(\sigma a)a t_X^{-1}$. A multiplicative $1$-cochain $\phi$ on $K$ assigns a
nonzero scalar $\phi(e)\in k^*$ to each edge $e$ of $K$, and its coboundary
on the triangle $(X,a)$ is
\[
 (d\phi)(X,a)=\phi(\sigma a)\phi(a)\phi(t_X)^{-1}.
\]
Rescaling each arrow $a$ of $\Gamma_A$ by $\phi(a)$ and dividing each mesh
relation $r_X$ by $\phi(t_X)$ multiplies the coefficient $c_{X,a}$ by
$(d\phi)(X,a)$ without changing the mesh ideal. Since $K$ has no $3$-cells
and is homotopy equivalent to a graph by Lemma~\ref{lem:mesh-walk-free}, we
have $H^2(K;k^*)=1$, so there exists a $1$-cochain $\phi$ with
$d\phi=c^{-1}$. This rescaling normalizes all mesh coefficients to $1$ and
induces an object-fixing isomorphism on $k(\Gamma_A)$.
\end{proof}

\subsection{The universal cover and the lifted mesh category}
\label{subsec:universal-cover}

By Proposition~\ref{prop:graded-realization}(3), we identify $\C$ with the
scalar mesh category $k(\Gamma_A)$ equipped with the coefficients
\eqref{eq:scalar-mesh}. Let $\widetilde K$ be the universal cover of $K$,
whose vertices are homotopy classes of walks in $K$ starting at a fixed
basepoint, and on which the fundamental group $G=\pi_1(K)$ acts freely by
deck transformations $\widetilde Y\mapsto g\widetilde Y$. Lifting the arrows
and translation edges of $\Gamma$ to $\widetilde K$ produces a translation
quiver $\widetilde\Gamma$. Because each two-arrow path $(\sigma a)a$ in a
mesh at $X$ is homotopic in $K$ to the translation edge $t_X$, all terms of
a mesh relation $r_X$ lift to paths in $\widetilde\Gamma$ sharing the same
source $\tau\widetilde X$ and target $\widetilde X$. Imposing the lifted
relations $r_{\widetilde X}$ with the same scalar coefficients $c_{X,a}$
defines the lifted scalar mesh category $\widetilde\C$.

Any \emph{contextual mesh relation}—that is, a relation of the form
$u r_Y v$ with $u$ and $v$ composable paths in $\Gamma$—lifts uniquely to
$\widetilde\Gamma$ once a lift of its starting vertex is chosen, and all of
its terms share the same lifted endpoints. Consequently, both the path space
and the mesh relation ideal decompose as direct sums over lifted endpoints,
so pushing paths down to $\Gamma$ induces a canonical linear isomorphism for
any chosen lifts $\widetilde X$ of $X$ and $\widetilde Y$ of $Y$:
\begin{equation}
 \bigoplus_{g\in G}\widetilde\C(\widetilde X,g\widetilde Y)
                 \xrightarrow{\sim}\C(X,Y).
 \label{eq:walk-hom-sum}
\end{equation}
In particular, the uniform bound on the length of nonzero paths in the
finite-dimensional category $\C$ also holds in $\widetilde\C$. Note also
that the lifted quiver $\widetilde\Gamma$ is connected: in a module category
every almost split sequence has a nonzero middle term, so every translation
edge $t_X$ in a walk on $\widetilde K$ can be replaced by a two-arrow path
$(\sigma a)a$ in $\widetilde\Gamma$.

Assigning length $1$ to each quiver arrow and length $2$ to each translation
edge $t_X$ preserves length across every triangle cell of $K$. Because closed
walks in the simply connected cover $\widetilde K$ are null-homotopic, this
edge length integrates on $\widetilde K$ to an integer-valued \emph{height
function} on the vertices of $\widetilde\Gamma$ that increases by $1$ along
every arrow. In particular, $\widetilde\Gamma$ is directed, every object of
$\widetilde\C$ has trivial endomorphism ring $k$, and only finitely many
vertices can be reached by directed paths between any two vertices of
prescribed heights.

\subsection{Finite subcategories and evaluation of lifted objects}
\label{subsec:lifted-projectives}

Let $\mathcal L$ denote the opposite of the full subcategory of
$\widetilde\C$ on the lifted projective vertices. Taking the opposite
reverses morphism directions so that they match those of
$A_0(x,y)=\C(P_y,P_x)$: choosing one reference lift $\widetilde x$ for each
projective vertex $P_x$, the covering decomposition \eqref{eq:walk-hom-sum}
gives a $G$-grading on $A_0$,
\[
 A_0(x,y)=\bigoplus_{g\in G}A_0(x,y)_g,\qquad
 A_0(x,y)_g=\mathcal L(\widetilde x,g\widetilde y)=\widetilde\C(P_{g\widetilde y},P_{\widetilde x}),
\]
which we refer to as the \emph{deck grading} to distinguish it from the
categorical integer grading of Section~\ref{sec:graded-category}. Since each
homogeneous component $A_0(x,y)_d$ of the categorical grading is at most
one-dimensional by Lemma~\ref{lem:ray-degree} and the two gradings are
compatible, each nonzero homogeneous ray line $A_0(x,y)_d$ is concentrated in
a single deck degree $g\in G$ and therefore lifts uniquely to a nonzero
morphism space $\mathcal L(\widetilde x,g\widetilde y)$ once the source lift
$\widetilde x$ is fixed.

More generally, both $\mathcal L$ and the finite categories constructed from
integer weights in Section~\ref{sec:local-ray-geometry} are instances of the
following standard covering construction for group-graded algebras.

\begin{definition}
\label{def:graded-covering-category}
Let $B$ be an algebra graded by a group $\mathcal G$, with a complete set of
primitive orthogonal idempotents $\{e_x\}$ homogeneous of identity degree.
The \emph{graded covering category} of $B$ has objects $(x,g)$ with
$g\in\mathcal G$, morphism spaces
\[
 \Hom\bigl((x,g),(y,h)\bigr)=B(x,y)_{g^{-1}h},
\]
and composition induced by the multiplication of $B$ (for an additively
written group $\mathcal G$, the degree $g^{-1}h$ is written $h-g$).
\end{definition}

Because the deck group $G$ acts freely on $\widetilde K$ and the chosen
lifts $\widetilde x$ form a transversal for the $G$-orbits of lifted
projectives, sending $g\widetilde x$ to $(x,g)$ identifies $\mathcal L$ with
the graded covering category of $A_0$ for its deck grading:
\[
 \mathcal L(g\widetilde x,h\widetilde y)
 =\mathcal L(\widetilde x,g^{-1}h\widetilde y)
 =A_0(x,y)_{g^{-1}h}.
\]

To transfer finite representation type from $A_0$ to finite full
subcategories of $\mathcal L$ (and later to subquotients in
Sections~\ref{sec:local-ray-geometry} and~\ref{sec:integral-partial-lifts}),
we use the standard induction--restriction argument for split bimodule
retractions (compare the cleaving-functor method in
\cite[Sections~3.2 and~3.8]{BGRS}).

\begin{lemma}
\label{lem:split-representation-transfer}
Let $C\subseteq B$ be a unital subalgebra of a finite-dimensional algebra $B$
over any field, let $q:C\twoheadrightarrow D$ be a quotient algebra, and
suppose that $q$ extends to a $C$-bimodule homomorphism $\pi:B\to D$. If $B$
is representation-finite, then so is $D$.
\end{lemma}
\begin{proof}
View any right $D$-module $M$ as a right $C$-module via $q$. The canonical
$C$-module homomorphism $M\to(M\otimes_C B)|_C$ sending $m$ to $m\otimes 1$
admits the $C$-linear retraction $m\otimes b\mapsto m\pi(b)$. Thus $M$ is a
direct summand of the restriction to $C$ of the $B$-module $M\otimes_C B$.
Since $B$ has only finitely many indecomposable modules up to isomorphism,
only finitely many indecomposable $C$-modules can arise as direct summands of
their restrictions, and since inflation along $q:C\twoheadrightarrow D$ is
fully faithful, $D$ is representation-finite.
\end{proof}

\begin{lemma}
\label{lem:homogeneous-projection}
Let $B$ be a representation-finite algebra graded by a group $\mathcal G$,
with primitive idempotents in the identity degree. Then every finite full
subcategory of the graded covering category of $B$ is representation-finite.
\end{lemma}
\begin{proof}
Let $H$ be the full subcategory on finitely many objects $(x_1,g_1),\ldots,(x_m,g_m)$,
so that $H(i,j)=B(x_i,x_j)_{g_i^{-1}g_j}$, and embed $H$ into the ungraded
category $B'$ on the same object indices with morphism spaces
$B'(i,j)=B(x_i,x_j)$. Because $B'$ is Morita equivalent to the idempotent
corner of $B$ supported on the distinct vertices among $x_1,\ldots,x_m$, it
is representation-finite by Lemma~\ref{lem:endpoint-chains}(1). Moreover, the
linear projection $\pi:B'\to H$ that projects each corner $B(x_i,x_j)$ onto
its homogeneous component of degree $g_i^{-1}g_j$ is an $H$-bimodule
retraction: for $p\in H(i,j)$, $b\in B'(j,l)$, and $q\in H(l,r)$, the product
with the homogeneous component of $b$ of degree $g_j^{-1}g_l$ is the unique
component of $pbq$ of degree $(g_i^{-1}g_j)(g_j^{-1}g_l)(g_l^{-1}g_r)=g_i^{-1}g_r$,
so $\pi(pbq)=p\pi(b)q$. Applying Lemma~\ref{lem:split-representation-transfer}
with $C=D=H$ shows that $H$ is representation-finite.
\end{proof}

\begin{lemma}
\label{lem:lifted-projective-evaluation}
The category $\mathcal L$ is connected and thin (meaning that
$\dim_k\mathcal L(u,v)\leq1$ for all objects $u,v$). For each object
$\widetilde x$ of $\mathcal L$, write $P_{\widetilde x}$ for the
corresponding lifted projective vertex in $\widetilde\C$, and for each object
$\widetilde X\in\widetilde\C$ define the $\mathcal L$-module
$V_{\widetilde X}$ by
\[
 V_{\widetilde X}(\widetilde x)=\widetilde\C(P_{\widetilde x},\widetilde X).
\]
Then each $V_{\widetilde X}$ is a finite-dimensional indecomposable
$\mathcal L$-module, and the evaluation functor
$\widetilde X\mapsto V_{\widetilde X}$ on $\widetilde\C$ is fully faithful.
\end{lemma}
\begin{proof}
Because the negative of height strictly increases along nonidentity morphisms
of $\mathcal L$, every object of $\mathcal L$ has endomorphism ring $k$, and
in any two-object full subcategory $\{u,v\}$ of $\mathcal L$ with
$\mathcal L(u,v)\ne0$ we have $\mathcal L(v,u)=0$. Every linear subspace of
$\mathcal L(u,v)$ is then a two-sided ideal of the two-object subcategory.
Since this subcategory is representation-finite by
Lemma~\ref{lem:homogeneous-projection}, it has only finitely many ideals by
Lemma~\ref{lem:endpoint-chains}(1), which forces $\dim_k\mathcal L(u,v)\leq1$.
Thus $\mathcal L$ is thin.

For any $\widetilde X\in\widetilde\C$, the uniform bound on nonzero path
lengths in $\widetilde\C$ implies that $V_{\widetilde X}$ is supported on
finitely many objects of $\mathcal L$. If $X\in\C$ denotes the projection of
$\widetilde X$, the covering isomorphism \eqref{eq:walk-hom-sum} gives, for
each reference projective lift $\widetilde x$,
\[
 \C(P_x,X)=\bigoplus_{g\in G}\widetilde\C(P_{g\widetilde x},\widetilde X)
 =\bigoplus_{g\in G}V_{\widetilde X}(g\widetilde x).
\]
Summing over $x$ identifies the evaluated $A_0$-module $\operatorname{ev}(X)$
of Proposition~\ref{prop:projective-evaluation} with the direct sum of the
spaces $V_{\widetilde X}(\widetilde z)$ over all objects
$\widetilde z\in\mathcal L$. Any homomorphism
$\varphi:V_{\widetilde X}\to V_{\widetilde Y}$ of $\mathcal L$-modules
acts componentwise on these spaces and commutes with the action of each
deck-homogeneous element of $A_0$, so it defines an $A_0$-module
homomorphism $\operatorname{ev}(X)\to\operatorname{ev}(Y)$ preserving each
individual summand indexed by $\widetilde z$. By
Proposition~\ref{prop:projective-evaluation}(1), this $A_0$-module map has
the form $\operatorname{ev}(f)$ for a unique $f\in\C(X,Y)$. Using
\eqref{eq:walk-hom-sum}, decompose $f=\sum_{g\in G}\bar f_g$, where
$\bar f_g\in\C(X,Y)$ is the projection of $f_g\in\widetilde\C(\widetilde X,g\widetilde Y)$.
Because composition with $f_g$ maps $V_{\widetilde X}(\widetilde z)$ into
$\widetilde\C(P_{\widetilde z},g\widetilde Y)=V_{\widetilde Y}(g^{-1}\widetilde z)$
and $G$ acts freely, the component $\operatorname{ev}(\bar f_g)$ for $g\ne1$
maps the summand at $\widetilde z$ into the distinct summand at
$g^{-1}\widetilde z$. Since $\varphi=\operatorname{ev}(f)$ preserves each
summand $\widetilde z$, we must have $\operatorname{ev}(\bar f_g)=0$ for all
$g\ne1$; faithfulness of $\operatorname{ev}$ and \eqref{eq:walk-hom-sum}
then give $f_g=0$ for $g\ne1$, so $f=\bar f_1$ and $\varphi$ is induced by
$f_1\in\widetilde\C(\widetilde X,\widetilde Y)$. This proves fullness, and
the same decomposition shows faithfulness. In particular,
$\End_{\mathcal L}(V_{\widetilde X})\cong\widetilde\C(\widetilde X,\widetilde X)=k$,
so each $V_{\widetilde X}$ is indecomposable.

Finally, since each $V_{\widetilde X}$ is indecomposable, its support lies in
a single connected component of $\mathcal L$. For every arrow
$\widetilde X\to\widetilde Y$ of $\widetilde\Gamma$, the induced nonzero
homomorphism $V_{\widetilde X}\to V_{\widetilde Y}$ forces the supports of
$V_{\widetilde X}$ and $V_{\widetilde Y}$ to lie in the same connected
component of $\mathcal L$. Because $\widetilde\Gamma$ is connected and every
object $\widetilde x\in\mathcal L$ belongs to the support of
$V_{P_{\widetilde x}}$, it follows that $\mathcal L$ is connected.
\end{proof}

\subsection{Extending a basis by one object}

We will construct a multiplicative basis of $\mathcal L$ inductively along an
exhaustion by finite full subcategories, adding one source or sink at a time.
When we adjoin a source $s$ to a finite directed thin category, the objects
$y$ receiving a nonzero morphism from $s$ form a finite poset, and
normalizing the new products amounts to trivializing a multiplicative
$1$-cocycle on the incidence category of that poset (compare Bretscher and
Gabriel \cite[Section~2.6]{BretscherGabriel}). Recall that the
\emph{incidence category} of a finite poset $R$ has the elements of $R$ as
objects and a single basis morphism $x\to y$ whenever $x\leq y$, with the
composite of any two comparable morphisms equal to the basis morphism between
their endpoints; an \emph{additive weight} on $R$ with values in an abelian
group $M$ is a function $w(x,y)\in M$ on comparable pairs $x\leq y$
satisfying $w(x,y)+w(y,z)=w(x,z)$ for $x\leq y\leq z$.

\begin{lemma}
\label{lem:poset-weights}
If the incidence algebra of a finite poset $R$ is representation-finite, then
every additive weight $w$ on $R$ with values in an abelian group $M$ has the
form $w(x,y)=\zeta_y-\zeta_x$ for some family of object values
$\zeta_x\in M$ (written multiplicatively as $w(x,y)=\zeta_y/\zeta_x$ when
$M=k^*$).
\end{lemma}
\begin{proof}
We proceed by induction on $|R|$. Choose a minimal element $s\in R$ and set
$U=\{y\in R:s<y\}\subseteq R\setminus\{s\}$. We claim that distinct
connected components of the subposet $U$ lie in distinct connected components
of $R\setminus\{s\}$. Suppose to the contrary that two connected components
$U_1,U_2$ of $U$ lie in the same connected component $C$ of
$R\setminus\{s\}$. For each scalar $\lambda\in k^*$, define a representation
of $R$ of dimension vector $(1,\ldots,1)$ by placing $k$ at every vertex, the
identity map $1_k$ on every relation inside $R\setminus\{s\}$ and from $s$ to
$U\setminus U_1$, and multiplication by $\lambda$ on every relation from $s$
to $U_1$. Any isomorphism between two such representations must act by a
constant scalar on the connected component $C$ and therefore cannot change
the ratio of the maps from $s$ into $U_1$ and $U_2$. Varying $\lambda\in k^*$
thus yields infinitely many pairwise nonisomorphic indecomposable
representations, contradicting the hypothesis that $R$ is
representation-finite.

Because the incidence algebra of $R\setminus\{s\}$ is an idempotent corner of
that of $R$, it is representation-finite by Lemma~\ref{lem:endpoint-chains}(1).
By the induction hypothesis, after subtracting a coboundary $\zeta_y-\zeta_x$
on $R\setminus\{s\}$, we may assume that $w(x,y)=0$ for all $x\leq y$ in
$R\setminus\{s\}$. For any $y\leq z$ in $U$, the additivity condition
$w(s,y)+w(y,z)=w(s,z)$ then gives $w(s,y)=w(s,z)$, so $w(s,-)$ is constant on
each connected component of $U$. Since distinct connected components of $U$
lie in distinct connected components of $R\setminus\{s\}$, we may add a
constant to $\zeta$ on each connected component of $R\setminus\{s\}$ to make
$w(s,y)=\zeta_y-\zeta_s$ for all $y\in U$ without altering the zero weights
on $R\setminus\{s\}$.
\end{proof}

To apply Lemma~\ref{lem:poset-weights} when adjoining a source $s$, we note
that the functor $y\mapsto T(s,y)$ on a finite directed thin category $T$
defines a thin module whose faithful support is a poset. We will reuse this
observation in Sections~\ref{sec:weights-and-parameters}
and~\ref{sec:local-ray-geometry}.

\begin{lemma}
\label{lem:thin-module-incidence}
Let $T$ be a finite directed category, and let $M$ be a covariant or
contravariant $k$-linear functor from $T$ to vector spaces satisfying
$\dim_k M(y)\leq1$ for every object $y$. Let $R$ be the category obtained
from $T$ by passing to the quotient by the annihilator ideal $\ann_T(M)$ and
discarding the objects $y$ with $M(y)=0$. Then $R$ is isomorphic to the
incidence category of a finite poset, and if $T$ is representation-finite,
then so is $R$.
\end{lemma}
\begin{proof}
Fix a nonzero basis vector in each one-dimensional space $M(y)$ for $y\in R$.
Because $R$ acts faithfully on $M$, every morphism space $R(y,z)$ acts
faithfully on the one-dimensional spaces $M(y)$ and $M(z)$ and therefore has
dimension at most one, spanned when nonzero by the unique morphism sending
the chosen basis vector to the chosen basis vector. Since $T$ is directed,
the relation $R(y,z)\ne0$ defines a partial order on the objects of $R$, and
these normalized morphisms identify $R$ with the incidence category of that
poset. Finally, if $T$ is representation-finite, then the full subcategory of
$T$ on the objects of $R$ is an idempotent corner of $T$ and hence
representation-finite by Lemma~\ref{lem:endpoint-chains}(1); since $R$ is a
quotient of that corner, $R$ is representation-finite as well.
\end{proof}

\begin{lemma}
\label{lem:source-basis-extension}
Let $T$ be a finite directed thin representation-finite category, let $s$ be
a source of $T$, and set $T'=T\setminus\{s\}$ and $M=\rad T(s,-)$. Let $R$ be
the representation-finite poset incidence category obtained from $T'$ and $M$
as in Lemma~\ref{lem:thin-module-incidence}.
\begin{enumerate}[label={\upshape(\arabic*)}]
\item Every multiplicative basis of $T'$ extends to a multiplicative basis of
$T$.
\item Let $\operatorname{Scal}(T)$ denote the abelian group of multiplicative
$k^*$-valued weights on the nonzero nonidentity morphism lines of $T$, that
is, assignments $d(r)\in k^*$ satisfying $d(rr')=d(r)d(r')$ whenever
$rr'\ne0$. Then restriction to $T'$ is a split surjection of abelian groups
\[
 \operatorname{Scal}(T)\longrightarrow \operatorname{Scal}(T')
\]
whose kernel is isomorphic to $(k^*)^{\pi_0(R)}$.
\end{enumerate}
Both assertions hold dually when $s$ is a sink of $T$.
\end{lemma}
\begin{proof}
Fix a multiplicative basis $\{b_{yz}\}$ of $T'$, and choose an arbitrary
nonzero vector $f_y\in T(s,y)$ for each object $y$ of $R$. Whenever a basis
element $b_{yz}\in T'(y,z)$ acts nontrivially on $M$, so that
$T(s,y)b_{yz}=T(s,z)\ne0$, we have $f_y b_{yz}=\mu_{yz}f_z$ for a nonzero
scalar $\mu_{yz}\in k^*$. Associativity of composition in $T$ implies that
$\mu_{yz}\mu_{zw}=\mu_{yw}$ for all comparable triples $y\leq z\leq w$ in
$R$. By Lemma~\ref{lem:poset-weights}, there exist scalars $\zeta_y\in k^*$
such that $\mu_{yz}=\zeta_z/\zeta_y$. Replacing each $f_y$ by $\zeta_y f_y$
gives $(\zeta_y f_y)b_{yz}=\zeta_z f_z$, extending the multiplicative basis
of $T'$ to $T$. This proves (1).

For (2), an element $d\in\operatorname{Scal}(T')$ extends to $d\in\operatorname{Scal}(T)$
if and only if its values $d(sy)\in k^*$ at the objects $y\in R$ satisfy
$d(sy)d(b_{yz})=d(sz)$ whenever $y\leq z$ in $R$. Choosing a base object
$y_D$ in each connected component $D$ of $R$ and a spanning tree of walks in
$D$ from $y_D$ to every $y\in D$, Lemma~\ref{lem:poset-weights} shows that we
obtain a well-defined homomorphic section
$\operatorname{Scal}(T')\to\operatorname{Scal}(T)$ by setting $d(sy_D)=1$ and
propagating $d(sy)$ along the walks in $D$ via products and inverses of the
values $d(b_{yz})$. Any two extensions of the same weight on $T'$ differ by a
constant scalar on each connected component $D\in\pi_0(R)$, so the kernel of
the restriction map is $(k^*)^{\pi_0(R)}$.
\end{proof}

We now apply Lemma~\ref{lem:source-basis-extension} along an exhaustion of
$\mathcal L$. Recall that a full subcategory of $\mathcal L$ is \emph{convex}
if it contains every object lying on a chain of nonzero morphisms between any
two of its objects.

\begin{lemma}
\label{lem:cover-exhaustion}
There exists an increasing sequence $\mathcal L_1\subset\mathcal L_2\subset\cdots$
of finite convex full subcategories of $\mathcal L$ whose union is
$\mathcal L$, such that each $\mathcal L_{m+1}$ is obtained from
$\mathcal L_m$ by adjoining a single source or sink. Every finite full
subcategory of $\mathcal L$ is directed, thin, and representation-finite.
Consequently, $\mathcal L$ admits a multiplicative basis, and the group
$\operatorname{Scal}(\mathcal L)$ of multiplicative $k^*$-valued weights on
$\mathcal L$ is an abelian divisible group.
\end{lemma}
\begin{proof}
Because $\mathcal L$ is connected and locally finite, it has countably many
objects, and because the integer height strictly decreases along every
nonidentity morphism of $\mathcal L$, there are only finitely many directed
paths between any two objects. The convex hull in $\mathcal L$ of any finite
set of objects is therefore finite. Enumerating the objects of $\mathcal L$
and taking the convex hulls of initial segments yields a cofinal sequence of
finite convex subcategories. Whenever $C\subsetneq C'$ is a proper inclusion
of finite convex subcategories, reachability along nonzero morphisms defines
a partial order on $C'$. If $C'\setminus C$ contains an element lying below
some element of $C$, any maximal such element of $C'\setminus C$ is a source
in the union of $C$ with that element, and adjoining it preserves convexity;
otherwise, any minimal element of $C'\setminus C$ retains convexity and
adjoins a sink. Refining the cofinal sequence in this way produces the
desired one-object exhaustion $\mathcal L_1\subset\mathcal L_2\subset\cdots$.

Every finite full subcategory of $\mathcal L$ is directed by height, thin by
Lemma~\ref{lem:lifted-projective-evaluation}, and representation-finite by
Lemma~\ref{lem:homogeneous-projection}. Applying
Lemma~\ref{lem:source-basis-extension}(1) at each step $\mathcal L_m\subset\mathcal L_{m+1}$
therefore builds a compatible sequence of multiplicative bases whose union is
a multiplicative basis of $\mathcal L$. Likewise, by
Lemma~\ref{lem:source-basis-extension}(2), the inverse system of weight groups
$\operatorname{Scal}(\mathcal L_m)$ consists of split surjections whose
kernels are products of copies of the divisible group $k^*$ (since $k$ is
algebraically closed). Hence the inverse limit
$\operatorname{Scal}(\mathcal L)=\varprojlim\operatorname{Scal}(\mathcal L_m)$
is isomorphic to a direct product of copies of $k^*$ and is therefore a
divisible abelian group.
\end{proof}

\subsection{Scalar automorphisms and descent}
\label{subsec:scalar-descent}

To make the multiplicative basis of $\mathcal L$ invariant under the action
of the free deck group $G=\pi_1(K)$, we first show that every multiplicative
scalar weight on $\mathcal L$ is a coboundary of object values.

\begin{lemma}
\label{lem:scalar-object-values}
Every multiplicative scalar weight $d\in\operatorname{Scal}(\mathcal L)$ is a
ratio of object values: there exist nonzero scalars $\zeta_x\in k^*$ for
$x\in\operatorname{ob}\mathcal L$ such that $d(r)=\zeta_y/\zeta_x$ for every
nonzero morphism line $r\in\mathcal L(x,y)$.
\end{lemma}
\begin{proof}
Each weight $d\in\operatorname{Scal}(\mathcal L)$ acts on $\mathcal L$ by an
object-fixing automorphism $\alpha_d$ that scales each nonzero morphism line
$r$ by $d(r)$. Precomposing $\mathcal L$-module actions with $\alpha_d^{-1}$
defines an action of the group $\operatorname{Scal}(\mathcal L)$ on the set
of isomorphism classes of finite-dimensional indecomposable
$\mathcal L$-modules. For any fixed finite subset
$S\subset\operatorname{ob}\mathcal L$, the full subcategory of $\mathcal L$
on $S$ is representation-finite by Lemma~\ref{lem:cover-exhaustion}, so there
are only finitely many isomorphism classes of indecomposable
$\mathcal L$-modules supported in $S$. Because $\operatorname{Scal}(\mathcal L)$
is a divisible abelian group by Lemma~\ref{lem:cover-exhaustion}, it has no
nontrivial finite quotient groups and must therefore act trivially on this
finite set of isomorphism classes. In particular, for every object
$\widetilde X\in\widetilde\C$, the evaluated module $V_{\widetilde X}$ of
Lemma~\ref{lem:lifted-projective-evaluation} satisfies
$V_{\widetilde X}\cong T(V_{\widetilde X})$, where $T$ denotes pullback along
$\alpha_d^{-1}$.

Choose an $\mathcal L$-module isomorphism
$\theta_{\widetilde X}:V_{\widetilde X}\xrightarrow{\sim}T(V_{\widetilde X})$
for each $\widetilde X\in\widetilde\C$, taking
$\theta_{P_{\widetilde x}}(v)=\alpha_d^{-1}(v)$ on the representable modules
$V_{P_{\widetilde x}}=\mathcal L(\widetilde x,-)$. By the full faithfulness
of $\widetilde X\mapsto V_{\widetilde X}$
(Lemma~\ref{lem:lifted-projective-evaluation}), transporting morphisms via
$f\mapsto\theta_{\widetilde Y}^{-1}\circ T(f)\circ\theta_{\widetilde X}$
defines an object-fixing automorphism $\Phi$ of $\widetilde\C$ whose
restriction to the opposite subcategory $\mathcal L$ of lifted projectives is
precisely $\alpha_d$.

Because the endpoints of each arrow $a:\widetilde U\to\widetilde V$ of
$\widetilde\Gamma$ differ in height by $1$, the space
$\widetilde\C(\widetilde U,\widetilde V)$ is one-dimensional and spanned by
$a$, so $\Phi(a)=w(a)a$ for some $w(a)\in k^*$. Since $\Phi$ preserves the
mesh relation $r_{\widetilde X}$ at each nonprojective vertex $\widetilde X$,
it scales $r_{\widetilde X}$ by some $\nu_{\widetilde X}\in k^*$, which
gives
\[
 w(\sigma a)w(a)=\nu_{\widetilde X}
\]
for every arrow $a$ ending at $\widetilde X$. Assigning $\nu_{\widetilde X}$
to each translation edge $t_{\widetilde X}$ makes $w$ a closed multiplicative
edge weight on the simply connected complex $\widetilde K$. Integrating $w$
along walks from a basepoint produces vertex scalars
$\zeta_{\widetilde X}\in k^*$ such that
$w(a)=\zeta_{\widetilde V}/\zeta_{\widetilde U}$ for every arrow
$a:\widetilde U\to\widetilde V$, and hence $\Phi$ acts on every morphism in
$\widetilde\C(\widetilde X,\widetilde Y)$ by multiplication by
$\zeta_{\widetilde Y}/\zeta_{\widetilde X}$. Since
$\mathcal L(\widetilde x,\widetilde y)=\widetilde\C(P_{\widetilde y},P_{\widetilde x})$
and $\Phi$ acts on $\mathcal L(\widetilde x,\widetilde y)$ by $d(\widetilde x,\widetilde y)$,
we obtain $d(\widetilde x,\widetilde y)=\zeta_{P_{\widetilde x}}/\zeta_{P_{\widetilde y}}$,
which has the required form on setting $\zeta_{\widetilde x}=\zeta_{P_{\widetilde x}}^{-1}$.
\end{proof}

We can now descend the multiplicative basis from $\mathcal L$ to $A_0$
(compare Bretscher and Gabriel \cite[Sections~3.1--3.2]{BretscherGabriel}).

\begin{proposition}
\label{prop:scalar-normalization}
Let $A_0$ be the standard form of $A$ constructed in
Construction~\ref{con:projective-algebra}, and let $P=\ray A_0$. Then there
is a graded algebra isomorphism $A_0\cong kP$ preserving the primitive
idempotents and the homogeneous ray lines. In particular, $A_0$ admits a
normed multiplicative basis.
\end{proposition}
\begin{proof}
Let $\{b_{xy}\}$ be the multiplicative basis of $\mathcal L$ provided by
Lemma~\ref{lem:cover-exhaustion}, where $b_{xy}\in\mathcal L(x,y)$ when
$\mathcal L(x,y)\ne0$. For each generator $g$ of the free deck group $G$, the
deck action on the basis defines a multiplicative weight $d_g\in\operatorname{Scal}(\mathcal L)$
by
\[
 g(b_{xy})=d_g(x,y)b_{gx,gy}.
\]
By Lemma~\ref{lem:scalar-object-values}, for each free generator $g$ there
exist scalars $a_g(x)\in k^*$ such that $d_g(x,y)=a_g(y)/a_g(x)$. Because
$G$ is a free group acting freely on $\operatorname{ob}\mathcal L$, each
$G$-orbit of objects is a Cayley tree for the free generators of $G$, so we
can solve the equations $f(gx)=a_g(x)f(x)$ simultaneously for all free
generators $g$ and all objects $x\in\operatorname{ob}\mathcal L$. Rescaling
the basis of $\mathcal L$ to
\[
 b'_{xy}=\frac{f(y)}{f(x)}b_{xy}
\]
preserves multiplicativity and makes the basis $G$-invariant:
$g(b'_{xy})=b'_{gx,gy}$.

As shown in Subsection~\ref{subsec:lifted-projectives}, each nonzero
homogeneous ray line of $A_0$ lies in a single deck degree and therefore
lifts uniquely to a nonzero morphism space $\mathcal L(x,y)$ once a lift $x$
of its source is chosen; by $G$-invariance, the image of $b'_{xy}$ in $A_0$
is independent of the choice of lift $x$. Since nonzero products of
homogeneous ray lines in $A_0$ lift to nonzero compositions in $\mathcal L$,
the descended vectors form a homogeneous multiplicative basis of $A_0$
realizing the isomorphism $A_0\cong kP$, and Corollary~\ref{cor:ray-transversal}
shows that this basis is normed.
\end{proof}

\begin{corollary}
\label{cor:graded-standard-form}
For $P=\ray A_0$, the linearization $kP\cong A_0$ is representation-finite
and satisfies $\add k(\Gamma_A)\simeq\modu kP$. Moreover, $A$ is standard in
the sense of Definition~\ref{def:standardness} if and only if $A\cong A_0$.
\end{corollary}
\begin{proof}
The first statement combines Propositions~\ref{prop:graded-realization}(3),
\ref{prop:projective-evaluation}, and~\ref{prop:scalar-normalization}. If $A$
is standard, then the object-preserving equivalence
$k(\Gamma_A)\simeq\ind(\modu A)$ restricts on the projective vertices to an
isomorphism $A_0\cong A$. Conversely, any algebra isomorphism $A\cong A_0$
induces an equivalence $F:\ind(\modu A)\xrightarrow{\sim}\ind(\modu A_0)\simeq\C$,
possibly permuting the vertices of $\Gamma_A$. Since $\C\cong k(\Gamma_A)$ is
generated in degree one with homogeneous relations, its associated graded
category $\gr_\JJ\C$ is canonically isomorphic to $\C$ itself, and the induced
graded equivalence $\gr_\JJ F:\C\xrightarrow{\sim}\C$ effects the same vertex
permutation as $F$. Composing $F$ with $(\gr_\JJ F)^{-1}$ therefore produces an
equivalence $\ind(\modu A)\xrightarrow{\sim}\C\cong k(\Gamma_A)$ fixing every
vertex.
\end{proof}
